\section*{Chapter \ref{ch:dv:the-structured-products-business} --- The Structured-Products Business}

\begin{solution}{exo:dv:the-structured-products-business:1}
A two-year zero-coupon bond of the issuer paying 100, and 0.6 at-the-money calls on 100 of the index expiring in two years. The
investor holds both: long the issuer's bond (and so its credit risk) and long the calls.
\end{solution}

\begin{solution}{exo:dv:the-structured-products-business:2}
$100e^{-0.06}=94.18$ at a zero spread and $100e^{-0.108}=89.76$ at 2.4\%: the spread adds 4.41 of option budget, more than doubling
the 4.32 left at a zero spread after the margin.
\end{solution}

\begin{solution}{exo:dv:the-structured-products-business:3}
The higher participation may be paid for by a weaker issuer's higher funding spread, which is the investor's compensation for a higher
risk of losing the protection in a default; or by a larger margin taken elsewhere. Compare notes at the same issuer credit, or add the
value of the credit risk to the comparison.
\end{solution}

\begin{solution}{exo:dv:the-structured-products-business:4}
The budget is $100-0.95\times94.18-1.5=9.03$, so the participation is $9.03/12.49=72.3\%$: five points of protection buy more than
twice the participation.
\end{solution}

\begin{solution}{exo:dv:the-structured-products-business:5}
The index's volatility is held near the target by construction, so the option is priced near the target volatility, well below the
underlying's, and its \omterm{def:dv:greeks-and-the-hedging-pnl:greeks}{vega} is small. The investor gives up exposure when volatility is high (after falls, typically) and takes the
rule's rebalancing, which sells after falls and buys after rallies; returns in trending markets differ from the plain index's.
\end{solution}

\begin{solution}{exo:dv:the-structured-products-business:6}
It lags by the excess of the decrement over the dividends, 2\% a year. If dividends fall to 1\%, the price index gains relative to the
total-return index while the \omterm{def:dv:the-structured-products-business:decrement}{decrement index} does not change its rule: the lag widens to 4\% a year, borne by the investor in the \omterm{def:dv:the-structured-products-business:decrement}{decrement
index}; the issuer, hedged on the fixed decrement, is unaffected.
\end{solution}

\begin{solution}{exo:dv:the-structured-products-business:7}
At a 1\% spread and a 1.5 margin: 5.17\% with the strike at 80\%, 6.96\% at 90\% and 9.71\% at 100\%. The funding benefit is the same, 3.92;
the puts sold are worth 2.55, 4.26 and 6.91 per 100.
\end{solution}

\begin{solution}{exo:dv:the-structured-products-business:8}
The participation depends at least as much on the issuer's funding spread (and on the margin) as on option prices: a weak issuer with the
same option prices offers more participation. The claim confuses credit with pricing.
\end{solution}

\begin{solution}{pb:dv:the-structured-products-business:1}
\textbf{1.} 20.5\%; 12.49 per 100.
\textbf{2.} $0.7\times12.49=8.74$.
\textbf{3.} The call is struck at the initial level, at the money; its price uses the smile's volatility at that strike and expiry.
\textbf{4.} Lower it, roughly as the volatility of the averaged level is lower (chapter 16), raising the participation.
\textbf{5.} Sell back the upside above the cap, lowering the option's price and raising the participation below the cap.
\textbf{6.} 94.18.
\textbf{7.} 34.6\%, 49.6\% and 64.2\%.
\textbf{8.} 2.40\% a year (240 basis points); the bond then costs 89.76.
\textbf{9.} The note is the issuer's debt: the guaranteed 100 is a promise of the issuer, valued at the rate the issuer pays to borrow.
\textbf{10.} The issuer's credit risk: in a default the protection, and the options' value, are lost with the bond.
\textbf{11.} 113.3\%: 100\% participation up to a 13.3\% gain.
\textbf{12.} A coupon of 6.96\% a year, with the capital exposed below 90\%.
\textbf{13.} About 62\% at a zero spread: at-the-money calls at the 10\% target volatility cost about 6.94 per 100 over two years.
\textbf{14.} It makes the note look generous (high participation) and removes the issuer's volatility and \omterm{def:dv:dividends-borrow-and-forwards:risk}{dividend risk}; the investor takes the
index's rule, and returns that can differ widely from the plain index's.
\textbf{15.} In the European Union: a key information document of at most three sides of A4, stand-alone, accurate, fair, clear and not
misleading, with standardised scenarios and costs.
\textbf{16.} It is fair if 240 basis points is the issuer's real credit spread and the client wants that credit risk; it is not a gift.
\textbf{17.} That the difference is the issuer's credit, and that the note should be compared with the issuer's bonds.
\textbf{18.} It lowers it: the bond leg is revalued at the wider spread, even if the index has not moved.
\textbf{19.} 2.40\% a year (240 basis points): the zero-coupon bond then costs 89.76, leaving 8.74 for 70\% of the 12.49 call after the 1.5 margin.
\textbf{20.} How much option the issuer's funding and the margin leave the investor, per unit of the underlying's rise.
\end{solution}

\begin{solution}{iq:dv:the-structured-products-business:1}
$100=\pi\,100e^{-(r+s)T}+p\,C+m$: the protected amount as the issuer's zero-coupon bond, the participation times the option leg's price, and
the margin. The participation rises with rates, with the issuer's spread and with lower protection, and falls with volatility, dividends
(through the call) and the margin.
\iqlookfor{the identity and each input's direction.}
\end{solution}

\begin{solution}{iq:dv:the-structured-products-business:2}
The bank receives the issue price now and owes the payoff later; hedging the payoff with derivatives leaves a loan at the bank's own funding
spread. Notes diversify the bank's funding sources and are priced off its funding curve.
\iqlookfor{the note as debt, and the hedged residual.}
\end{solution}

\begin{solution}{iq:dv:the-structured-products-business:3}
A note paying an enhanced coupon and repaying the underlying's value, rather than par, if it ends below a strike: the investor is short a put
(and long the issuer's bond). The coupon is the funding rate plus the put premium, less the margin.
\iqlookfor{the short put and where the coupon comes from.}
\end{solution}

\begin{solution}{iq:dv:the-structured-products-business:4}
The index's volatility is managed to the target, so options on it are priced near the target and have little \omterm{def:dv:greeks-and-the-hedging-pnl:greeks}{vega}. What remains for the issuer
is the gap between the rule's estimated and realised volatility (sudden jumps before the exposure adjusts) and the index's rebalancing costs;
the issuer also bears model risk on the index's path-dependence.
\iqlookfor{the managed volatility and the residual risks.}
\end{solution}

\begin{solution}{iq:dv:the-structured-products-business:5}
The bond leg of every outstanding note is worth less (discounted at the wider spread), so secondary prices fall; an issuer that carries its notes at
fair value shows a gain on its own debt. New issues can offer better terms. Funding plans change if buy-backs rise.
\iqlookfor{secondary prices, own-credit accounting and new issuance.}
\end{solution}

\begin{solution}{iq:dv:the-structured-products-business:6}
Sell more optionality: a put closer to the money (5.2\%, 7.0\%, 9.7\% for strikes of 80\%, 90\%, 100\% in the chapter), a daily barrier, a worst-of on
several names (correlation), an autocall that shortens the note; use a weaker issuer (credit); cut the margin. Each raises the coupon by selling the
investor more risk; explain each in those terms.
\iqlookfor{the levers and the risk each transfers.}
\end{solution}
